\chapter{Gaussian Information Geometry}
\label{ch:infogeometry}

\Cref{ch:expfamily} distinguishes general law fibers from smooth
statistical manifolds, regular exponential families, and the
multivariate-Gaussian realization. This chapter works only in the last
tier. It identifies the dual Gaussian charts, computes their Fisher
metrics, and separates frame-invariant generalized spectral data from
ordinary chart-dependent eigenvalues. It does not assert that an RG flow
contracts any information-geometric quantity.

\section{Canonical and expectation coordinates}
\label{sec:ig-expfam}

Use the nondegenerate Gaussian family and log normalizer of
\Cref{prop:gauss-log-normalizer-moments}. For the stacked latent
\(Y\in\R^n\), \(n=NK\), write
\begin{equation}
q(y)=\exp\left(h^\top y-\tfrac12 y^\top Jy-\mathsf A(h,J)\right),
\qquad h\in\R^n,\quad J\in\Sym_{++}^n,
\label{eq:ig-density}
\end{equation}
where
\begin{equation}
\mathsf A(h,J)
=\tfrac12h^\top J^{-1}h-\tfrac12\log\det J+\tfrac n2\log2\pi.
\label{eq:ig-lognormalizer}
\end{equation}
For sufficient statistic \(T(y)=(y,yy^\top)\), the canonical natural and
expectation coordinates are
\begin{align}
\eta&=\left(h,-\tfrac12J\right),
\label{eq:ig-natural}\\
\tau&=\nabla_\eta\mathsf A
=\left(\mu,M_2\right),
\qquad M_2=C+\mu\mu^\top=\E_q[YY^\top],
\label{eq:ig-expectation}
\end{align}
The information-form chart \((h,J)\) ranges over
\(\R^n\times\Sym_{++}^n\), whereas the canonical matrix coordinate
\(-J/2\) is negative definite. The family is regular and minimal, so the
Legendre map is a bijection on its natural domain
\citep{Brown1986}. \status{ESTABLISHED}

The Fisher matrix in natural coordinates is the Hessian of the potential and equals the covariance of the sufficient statistic,
\begin{equation}
g_\eta=\nabla^2_\eta\mathsf A(\eta)=\operatorname{Cov}_q\left(T(Y)\right),
\label{eq:ig-fisher-hessian}
\end{equation}
which is the standard exponential-family fact \citep{Brown1986,AmariNagaoka2000,Amari2016}. Because $\tau=\nabla_\eta\mathsf A$, the Jacobian of the coordinate change is $\partial\tau/\partial\eta=g_\eta$, and the Fisher metric expressed in the expectation coordinates is
\begin{equation}
g_\tau=\left(\frac{\partial\eta}{\partial\tau}\right)^{\top} g_\eta\left(\frac{\partial\eta}{\partial\tau}\right)=g_\eta^{-1},
\label{eq:ig-dual-inverse}
\end{equation}
the Hessian of the convex conjugate $\mathsf A^{*}(\tau)=\sup_\eta\left(\langle\eta,\tau\rangle-\mathsf A(\eta)\right)$, which for this family is the negative differential entropy of $q$ \citep{AmariNagaoka2000,nielsen2020elementary}. \status{ESTABLISHED}

Four coordinate presentations must be kept apart. The canonical natural
chart is \(\eta=(h,-J/2)\); the affinely equivalent information-form chart
is \((h,J)\); the expectation chart is
\(\tau=(\mu,M_2)\); and the moment chart is \((\mu,C)\). The last is useful
for sampling and covariance transport but is not Legendre dual to the
natural chart, because \(C=M_2-\mu\mu^\top\). \status{DEFINITION}

Symmetric-matrix coordinates carry one further convention. Derivatives in the matrix argument are taken with respect to the Frobenius pairing $\langle X,Z\rangle=\Tr(X^\top Z)$ on the space of symmetric matrices. Concretely this is the chart whose coordinates are the diagonal entries together with the off-diagonal entries scaled by $\sqrt2$, for which the Euclidean inner product of coordinate vectors equals the Frobenius pairing of the matrices they represent. With the unscaled half-vectorization instead, every off-diagonal coordinate derivative acquires a factor of two and the matrix representing \eqref{eq:ig-dual-inverse} changes accordingly. All matrix blocks displayed below are in the Frobenius-paired chart. \status{DEFINITION}

When the recognition precision has the interaction form of
\Cref{ch:gaussian}, \(J=\Lambda\), with
\(\Lambda_{ii}=A_i+\sum_{j\neq i}W_{ij}\),
\(\Lambda_{ij}=-W_{ij}\), and \(C=\Lambda^{-1}\). Below, \(\Lambda\)
emphasizes the interaction operator and \(J\) the natural parameter; they
are the same matrix.

\section{The chart correction}
\label{sec:ig-chart}

It is often asserted that the mean sector of the Fisher metric of a Gaussian family \emph{is} the precision. The assertion is true, and it is the standard statement of Gaussian information geometry, but it is true in the moment chart. It is false in the expectation chart, which is the chart the exponential-family duality of Section~\ref{sec:ig-expfam} actually delivers. Both facts are proved below, and the discrepancy between them is not a small correction.

\medskip
\propositionheading{Fisher metric in the moment chart}{prop:ig-fisher-moment-chart}
\emph{In the chart $(\mu,C)$ the Fisher metric of the Gaussian family is block diagonal, and for tangent vectors $(u,A)$ and $(v,B)$ with $u,v\in\R^{n}$ and $A,B$ symmetric,}
\begin{equation}
g_{(\mu,C)}\big((u,A),(v,B)\big)=u^\top\Lambda v+\tfrac12\Tr\left(\Lambda A\Lambda B\right),
\qquad
\Lambda=C^{-1}.
\label{eq:ig-moment-metric}
\end{equation}
\status{ESTABLISHED}

\noindent\emph{Proof.} Write $e=y-\mu$ and $\log q(y)=-\tfrac12 e^\top\Lambda e+\tfrac12\log\det\Lambda-\tfrac{n}{2}\log2\pi$. A displacement $(u,A)$ in the moment chart moves $\mu$ by $u$ and $C$ by $A$, hence $\Lambda$ by $-\Lambda A\Lambda$, and the directional score is
\begin{equation}
\partial_{(u,A)}\log q=u^\top\Lambda e+\tfrac12 e^\top\Lambda A\Lambda e-\tfrac12\Tr(\Lambda A).
\label{eq:ig-score}
\end{equation}
Under $e\sim\mathcal N(0,C)$ the first term has second moment $u^\top\Lambda C\Lambda u=u^\top\Lambda u$. The cross term between the linear and quadratic pieces is a third central moment of a Gaussian and vanishes, which is exactly why the metric is block diagonal in this chart. Setting $S=\Lambda A\Lambda$, the remaining piece is $\tfrac14\operatorname{Var}(e^\top Se)=\tfrac12\Tr(SCSC)=\tfrac12\Tr(\Lambda A\Lambda A)$. Polarization gives \eqref{eq:ig-moment-metric}. \hfill$\square$

\medskip
Expression \eqref{eq:ig-moment-metric} is the classical Fisher--Rao metric of the multivariate normal model; it is the metric whose geodesic equations and Rao distance are studied by \citet{Skovgaard1984} and by \citet{AtkinsonMitchell1981}, and it appears in the survey of \citet{PineleStrapassonCosta2020} as their equation~(7), in the general reparameterized form $g_{ij}=\partial_i\mu^\top\Sigma^{-1}\partial_j\mu+\tfrac12\Tr\left(\Sigma^{-1}\partial_i\Sigma \Sigma^{-1}\partial_j\Sigma\right)$. The printed equation in that reference carries a typographical slip, repeating the index $i$ in place of $j$ in the second factor of the trace; the intended and correct expression is the one displayed here, and it agrees with \eqref{eq:ig-moment-metric} on setting $\theta=(\mu,C)$. That reference also records, as its equation~(8), that the affine action $(\mu,\Sigma)\mapsto(Q\mu+c,Q\Sigma Q^\top)$ is an isometry of this metric for every $Q\in\GL_n(\R)$ and $c\in\R^{n}$, a fact used again in Section~\ref{sec:ig-gauge}.

\medskip
\propositionheading{Fisher metric in the expectation chart}{prop:ig-fisher-expectation-chart}
\emph{In the chart $(\mu,M_2)$ with $M_2=C+\mu\mu^\top$, the Fisher metric has mean block}
\begin{equation}
g_{[\mu\mu]}=\left(1+\mu^\top\Lambda\mu\right)\Lambda+\Lambda\mu\mu^\top\Lambda ,
\label{eq:ig-meanblock}
\end{equation}
\emph{cross block}
\begin{equation}
g_{[\mu M_2]}\big(u,B\big)=-u^\top\Lambda B\Lambda\mu ,
\label{eq:ig-crossblock}
\end{equation}
\emph{and second block $g_{[M_2M_2]}(A,B)=\tfrac12\Tr(\Lambda A\Lambda B)$, unchanged from \eqref{eq:ig-moment-metric}. Hence the metric is not block diagonal in this chart unless $\mu=0$.}
\status{ESTABLISHED}

\noindent\emph{Proof.} The chart change is $C=M_2-\mu\mu^\top$, so a displacement $(u,B)$ in the expectation chart induces
\begin{equation}
\delta\mu=u,
\qquad
\delta C=B-\left(u\mu^\top+\mu u^\top\right),
\label{eq:ig-chartjacobian}
\end{equation}
and the metric in the new chart is the pullback of \eqref{eq:ig-moment-metric} along \eqref{eq:ig-chartjacobian}. For the mean block take $B=0$, so $\delta C=-(u\mu^\top+\mu u^\top)$. Then
\begin{align}
\Tr\left(\Lambda \delta C \Lambda \delta C\right)
&=\Tr\left[\left(\Lambda u\mu^\top+\Lambda\mu u^\top\right)\left(\Lambda u\mu^\top+\Lambda\mu u^\top\right)\right]
\notag\\
&=2\left(u^\top\Lambda\mu\right)^2+2\left(\mu^\top\Lambda\mu\right)\left(u^\top\Lambda u\right),
\label{eq:ig-crossterm}
\end{align}
each of the four traces being evaluated by cycling. Adding the mean contribution $u^\top\Lambda u$ from \eqref{eq:ig-moment-metric} and halving \eqref{eq:ig-crossterm} gives the quadratic form
\begin{equation}
u^\top\Lambda u+\left(\mu^\top\Lambda\mu\right)\left(u^\top\Lambda u\right)+\left(u^\top\Lambda\mu\right)^2
=u^\top\left[\left(1+\mu^\top\Lambda\mu\right)\Lambda+\Lambda\mu\mu^\top\Lambda\right]u,
\label{eq:ig-meanquad}
\end{equation}
which is \eqref{eq:ig-meanblock}. For the cross block pair $(u,0)$ with $(0,B)$: the mean parts do not meet, and the covariance parts contribute $\tfrac12\Tr\left(\Lambda\left[-(u\mu^\top+\mu u^\top)\right]\Lambda B\right)=-u^\top\Lambda B\Lambda\mu$, which is \eqref{eq:ig-crossblock}. For the second block take $u=0$, so $\delta C=B$ and nothing changes. \hfill$\square$

\medskip
\corollaryheading{Size of the discrepancy and the exact equality case}{cor:ig-mean-block-discrepancy}
\emph{Write $v=\Lambda^{1/2}\mu$. Then}
\begin{equation}
g_{[\mu\mu]}=\Lambda^{1/2}\left[\left(1+\|v\|^2\right)I_{n}+vv^\top\right]\Lambda^{1/2},
\label{eq:ig-meanblock-factored}
\end{equation}
\emph{so the eigenvalues of $\Lambda^{-1}g_{[\mu\mu]}$ are $1+\mu^\top\Lambda\mu$ with multiplicity $n-1$ and $1+2\mu^\top\Lambda\mu$ with multiplicity one. Consequently $g_{[\mu\mu]}\succeq\Lambda$ always, and $g_{[\mu\mu]}=\Lambda$ if and only if $\mu=0$.}
\status{ESTABLISHED}

\noindent\emph{Proof.} Equation \eqref{eq:ig-meanblock-factored} is \eqref{eq:ig-meanblock} rewritten. The bracket is $\left(1+\|v\|^2\right)I+vv^\top$, whose eigenvalues are $1+\|v\|^2+\|v\|^2$ along $v$ and $1+\|v\|^2$ on $v^\perp$; since $\Lambda^{-1}g_{[\mu\mu]}=\Lambda^{-1/2}\left[\cdot\right]\Lambda^{1/2}$ is similar to the bracket, its spectrum is that list. For the equality case, $g_{[\mu\mu]}-\Lambda=\left(\mu^\top\Lambda\mu\right)\Lambda+\Lambda\mu\mu^\top\Lambda$ is a sum of two positive semidefinite matrices, so it vanishes only if both vanish; the first forces $\mu^\top\Lambda\mu=0$ and hence $\mu=0$ because $\Lambda\succ0$. The converse is immediate. \hfill$\square$

\medskip
The discrepancy is therefore governed by the squared Mahalanobis norm of the mean, $\mu^\top\Lambda\mu$, and it is unbounded. It is not a normalization constant and it does not disappear in any limit that keeps the mean away from the origin.

The mechanism behind the discrepancy is worth stating plainly, because it is what makes the error easy to commit. Both charts call their first coordinate $\mu$, but they do not have the same coordinate vector fields. The vector $\partial/\partial\mu$ at fixed $C$ moves the mean while holding the covariance fixed, so \(M_2\) changes; the vector $\partial/\partial\mu$ at fixed $M_2$ moves the mean and compensates in $C$. These are different tangent vectors at the same point of the manifold, and a metric assigns them different lengths. Writing $g_{[\mu\mu]}$ without saying which of the two is meant does not denote anything.

The same observation disposes of the phrase ``the mean sector'' in the expectation chart. A metric tensor has a well defined restriction to a subspace of the tangent space only once that subspace is declared; what is not chart independent is the \emph{complement}. In the moment chart the covariance directions are Fisher-orthogonal to the mean directions, by the vanishing third-moment argument in the proof of \Cref{prop:ig-fisher-moment-chart}, so ``the mean sector'' can be taken to mean either the restriction to the mean directions or the orthogonal quotient, and the two agree. In the expectation chart \eqref{eq:ig-crossblock} is nonzero, the mean and second-moment directions are not Fisher-orthogonal, and the restriction \eqref{eq:ig-meanblock} differs from the quotient obtained by projecting out the second-moment directions. Absent a declared splitting, ``the mean sector'' in that chart names no unique object.

\medskip
\noindent\textbf{The chart discipline.} A coordinate identity that mixes the mean and covariance sectors of a Gaussian family must name the chart in which it holds. Identity \eqref{eq:ig-moment-metric} holds in $(\mu,C)$; identities \eqref{eq:ig-meanblock} and \eqref{eq:ig-crossblock} hold in $(\mu,M_2)$; by \Cref{cor:ig-mean-block-discrepancy} the two mean blocks agree only at $\mu=0$. \status{ESTABLISHED}

\subsection*{Verification}

The statements are proved above in general dimension; the machine checks reported here test those proofs and are not a substitute for them. Two symbolic routes and two numerical routes were run, all in exact or double precision on the CPU.

The first symbolic route is independent of \Cref{prop:ig-fisher-moment-chart}: the Hessian of \eqref{eq:ig-lognormalizer} was formed by computer algebra in the natural chart $(h,\operatorname{svec}(-J/2))$ with the Frobenius-paired symmetric coordinates, evaluated at exact rational $(\mu,\Lambda)$, inverted exactly as in \eqref{eq:ig-dual-inverse}, and its blocks compared with \eqref{eq:ig-meanblock} and \eqref{eq:ig-crossblock}. At dimensions $n=1,2,3$ the difference was the exact zero matrix in both blocks, and at $n=2$ the cross block became exactly zero when $\mu=0$ was substituted and was nonzero otherwise, as \Cref{prop:ig-fisher-expectation-chart} requires. The second symbolic route carries general symbols: at dimensions $n=2$ and $n=3$, with $\mu$ a vector of free symbols and $\Lambda$ a general symmetric matrix of free symbols, the pullback of \eqref{eq:ig-moment-metric} along \eqref{eq:ig-chartjacobian} was simplified to zero difference against all three blocks of \Cref{prop:ig-fisher-expectation-chart}. Those two runs are proofs at those dimensions, not samples.

The deterministic numerical route \texttt{CHK-IG-EXPECTATION-METRIC} uses the closed-form Gaussian relative entropy and its Fisher quadratic limit in six directions at a fully specified base point. It compares the expectation-chart block with Equation~\eqref{eq:ig-meanblock}, uses the moment-chart precision as a negative control against chart conflation, and checks the predicted generalized spectrum. The computation is reproducible corroboration, not a proof. \status{NUMERICAL}

\section{Sample space and parameter space}
\label{sec:ig-typing}

The renormalization chapters will import vocabulary from a literature that computes Fisher metrics of likelihoods with respect to \emph{model} parameters and reads their spectra as a scale \citep{BermanKlingerStapleton2023}. The object developed above is not that object, and keeping the two typed apart is a precondition for importing anything at all.

Three spaces are in play. The latent sample space is $\R^{n}$, where $Y$ takes its values; the interaction energy $y\mapsto\tfrac12 y^\top\Lambda y$ is a quadratic form on it. The recognition parameter manifold is the set of laws \eqref{eq:ig-density}, coordinatized by any of the four charts of Section~\ref{sec:ig-expfam}; the Fisher metric \eqref{eq:ig-moment-metric} is a Riemannian metric on it. The model parameter space is where $\theta$ lives, indexing the generative kernel; the Fisher metric of the likelihood $\theta\mapsto p_\theta(o\given X)$ is a Riemannian metric on \emph{that}, and its dimension is the dimension of $\theta$, which has no relation to $n=NK$.

What makes the confusion easy is an exact but scoped component identity in this family, and stating it precisely is more useful than warning against it. The recognition family is a location family in its mean, so the mean parameter takes values in the same vector space $\R^{n}$ as the sample variable, and the mean block \eqref{eq:ig-moment-metric} of the parameter-space metric has the same components as the sample-space form $\Lambda$. The equality is of components in matched charts, not of objects: one is a metric on the tangent spaces of a manifold of probability laws, the other is a bilinear form on the latent vector space. For a family that is not a location family the two have different dimensions and no comparison arises. The identity is therefore a feature of the Gaussian mean sector and cannot be used as evidence that a construction defined on one space transfers to the other. \status{ESTABLISHED}

There is a corresponding exact but scoped relation to the interaction
Laplacian \(L\).  If transported Gaussian laws have covariance fixed within
the mean submodel and parallel along each edge, and the edge quadratic uses
that transported precision, their pairwise KL is exactly the associated
mean-sector Fisher quadratic; summing over edges gives the connection
Laplacian form.  For a regular non-Gaussian family, freezing the Fisher
tensor at one parallel base law gives the analogous form only as the local
second-order term of transported KL \citep{AmariNagaoka2000}.  Outside these
Gaussian or frozen-local hypotheses, a sample-space \(L\) is not a
recognition Fisher tensor or a natural-gradient operator.
\status{ESTABLISHED}

The distinction has a sharp consequence for coarse-graining, which the next chapters will need. Restricting the recognition mean to a subspace and marginalizing onto that subspace are different operations and give different metrics.

\medskip
\propositionheading{Pullback versus pushforward}{prop:ig-pullback-vs-pushforward}
\emph{Let $\Lambda\succ0$ on $\R^{n}$, let $B\in\R^{n\times k}$ have orthonormal columns, and let $B_\perp$ span the orthogonal complement. The Fisher metric of the submodel obtained by constraining the mean to $\operatorname{range}(B)$ is the restriction $B^\top\Lambda B$; the Fisher metric of the mean of the pushforward law under $y\mapsto B^\top y$ is the marginal precision $\left(B^\top\Lambda^{-1}B\right)^{-1}$. They satisfy}
\begin{equation}
B^\top\Lambda B-\left(B^\top\Lambda^{-1}B\right)^{-1}
=\left(B^\top\Lambda B_\perp\right)\left(B_\perp^\top\Lambda B_\perp\right)^{-1}\left(B_\perp^\top\Lambda B\right)
\succeq 0,
\label{eq:ig-restriction-marginal}
\end{equation}
\emph{with equality if and only if $B^\top\Lambda B_\perp=0$, that is if and only if the retained and discarded subspaces are $\Lambda$-orthogonal.}
\status{ESTABLISHED}

\noindent\emph{Proof.} The matrix $\left[B\quad B_\perp\right]$ is orthogonal, so in that basis $\Lambda$ has blocks $\Lambda_{11}=B^\top\Lambda B$, $\Lambda_{12}=B^\top\Lambda B_\perp$, $\Lambda_{22}=B_\perp^\top\Lambda B_\perp$. The block inversion formula gives $B^\top\Lambda^{-1}B=\left(\Lambda_{11}-\Lambda_{12}\Lambda_{22}^{-1}\Lambda_{21}\right)^{-1}$ \citep{HornJohnson2013}, and inverting again yields \eqref{eq:ig-restriction-marginal}. The right side is a congruence of $\Lambda_{22}^{-1}\succ0$ and is therefore positive semidefinite, and it vanishes exactly when $\Lambda_{12}=0$. The identification of the two sides as pullback and pushforward metrics is the Gaussian conditional and marginal precision respectively \citep{Lauritzen1996}. \hfill$\square$

\medskip
Restriction is a pullback along an injective map of parameters, and it makes the metric larger in the Loewner order; marginalization is a pushforward along a map of the sample space, and it makes the metric smaller. The monotonicity theorems of information geometry govern the second operation, not the first, a point taken up in Section~\ref{sec:ig-notclaimed}. The deterministic replacement check \texttt{CHK-IG-PULLBACK-PUSHFORWARD} tests the Schur identity and Loewner order over a specified ensemble and includes the $\Lambda$-orthogonal equality control. Computation is not proof and the proof is above. \status{NUMERICAL}

Whether the scale-setting construction of the model-space literature transfers to $\Lambda$ is therefore not settled by observing that both are called Fisher metrics. The transfer requires a declaration that makes the coarse recognition mean a parameter of the generative law fitted against the bound, after which the relevant Fisher metric is the one with respect to that parameter and the import becomes a statement rather than an analogy. This section records the requirement and does not discharge it. What would close it is that declaration together with a check that the resulting parameter-space metric is nondegenerate, since the construction in that literature is written throughout with an inverted metric and does not treat the singular case. \status{OPEN}

\section{Gauge invariants of the interaction precision}
\label{sec:ig-gauge}

A local reframing is a change of local trivialization of the fibers, and it changes the components of every object while changing no law. Let
\begin{equation}
T=\operatorname{diag}\left(T_1,\dots,T_N\right),
\qquad
T_i\in\GL^{+}(K),
\label{eq:ig-reframing}
\end{equation}
act on the assembled latent coordinate by $y\mapsto Ty$. Every quadratic form in $y$ then transforms by inverse congruence,
\begin{equation}
\Lambda\longmapsto T^{-\top}\Lambda T^{-1},
\qquad
L\longmapsto T^{-\top}LT^{-1},
\qquad
A\longmapsto T^{-\top}AT^{-1},
\label{eq:ig-congruence}
\end{equation}
where $L$ is the Laplacian part of the interaction precision and $A=\operatorname{diag}(A_1,\dots,A_N)$ collects the self terms, so that $\Lambda=L+A$. Two hypotheses are being used and should be visible. First, the splitting $\Lambda=L+A$ is a declared decomposition of the assembled precision, made once in the frame-dependent chart; \eqref{eq:ig-congruence} then says that congruence acts on the two summands separately, which is immediate from linearity. Second, the reframing is block diagonal with respect to the agent partition, which is what a change of local trivialization is; a coordinate change mixing agents would not preserve the block structure that makes $A$ block diagonal. \status{HYPOTHESIS}

Under a nonconstant reframing the image $T^{-\top}LT^{-1}$ need not itself be a matrix-weighted Laplacian, since the row sums that define that normal form are not preserved when neighboring blocks are reframed differently. This is a statement about normal forms, not about the object: $L$ is a bilinear form on the latent space, and \eqref{eq:ig-congruence} is how a bilinear form reads in a new chart. The invariants below compare forms and are therefore unaffected.

\medskip
\propositionheading{Frame invariance of the generalized spectrum}{prop:ig-generalized-spectrum-invariance}
\emph{Let $\Lambda\succ0$ and $L=L^\top$, and let $T$ be invertible. The generalized eigenvalue problem}
\begin{equation}
Lx=d\Lambda x
\label{eq:ig-pencil}
\end{equation}
\emph{and the reframed problem $\left(T^{-\top}LT^{-1}\right)x'=d'\left(T^{-\top}\Lambda T^{-1}\right)x'$ have the same eigenvalues, with eigenvectors related by $x'=Tx$.}
\status{ESTABLISHED}

\noindent\emph{Proof.} Suppose $Lx=d\Lambda x$ and set $x'=Tx$. Then $T^{-\top}LT^{-1}x'=T^{-\top}Lx=d T^{-\top}\Lambda x=d\left(T^{-\top}\Lambda T^{-1}\right)x'$. The map $x\mapsto Tx$ is a bijection, so the correspondence is onto, and the argument run with $T^{-1}$ gives the converse. \hfill$\square$

\medskip
Because $\Lambda\succ0$, the pair $(L,\Lambda)$ is a symmetric-definite, hence regular, pencil: it is simultaneously diagonalizable by congruence, all $n$ generalized eigenvalues are real, and they obey a Rayleigh characterization \citep{HornJohnson2013,GolubVanLoan2013}. The positive-definiteness hypothesis is load bearing. Congruence covariance alone also holds for a singular pencil, but it does not imply that the determinant polynomial has degree $n$ or even that it is nonzero. For example, at a common singular ray $(L,\Lambda)=(L,L)$ with $\ker L\neq\{0\}$,
\[
\det(L-d\Lambda)=\det\bigl((1-d)L\bigr)\equiv0,
\]
so the full-space determinant does not define an $n$-element generalized spectrum. A quotient, pin, or mass repair must first be declared and the resulting pencil proved regular before the conclusions below can be invoked. Thus \Cref{prop:ig-generalized-spectrum-invariance,prop:ig-generalized-spectrum-localization,prop:ig-frame-dependent-spectra-determinants} apply to the Gaussian probability interior $\Lambda\succ0$ (or to a separately proved regular repair), not automatically to an operator-layer boundary point. That qualification is the pencil-level instance of the hierarchy in \Cref{ch:expfamily}.

\medskip
\propositionheading{Localization and endpoints}{prop:ig-generalized-spectrum-localization}
\emph{If in addition $L\succeq0$ and $A=\Lambda-L\succeq0$, then every generalized eigenvalue of \eqref{eq:ig-pencil} lies in $[0,1]$. The value $0$ is attained exactly on $\ker L$, whose dimension is at least $K$ because every configuration constant across agents lies in it, and the value $1$ exactly on $\ker A$.}
\status{ESTABLISHED}

\noindent\emph{Proof.} By the Rayleigh characterization, each generalized eigenvalue is a value of $d(x)=x^\top Lx/x^\top\Lambda x$ at a stationary point, and $x^\top\Lambda x=x^\top Lx+x^\top Ax$ with both terms nonnegative and the sum positive. Hence $0\le d(x)\le1$ for all $x\neq0$, with $d(x)=0$ iff $x^\top Lx=0$, which for $L\succeq0$ is $x\in\ker L$, and $d(x)=1$ iff $x\in\ker A$. That the constant configurations $\mathbf 1\otimes w$, $w\in\R^{K}$, lie in $\ker L$ follows from the difference form $x^\top Lx=\sum_{i<j}\left(x_i-x_j\right)^\top W_{ij}\left(x_i-x_j\right)$. \hfill$\square$

\medskip
The generalized eigenvalue is therefore readable, and its reading is frame independent: $d_k$ is the fraction of the total precision along the $k$-th pencil direction that is supplied by interaction rather than by self terms. The consensus directions sit at $d=0$ because interaction supplies no precision there; a direction that no self term constrains sits at $d=1$.

\medskip
\propositionheading{Frame dependence of absolute spectra and subspace determinants}{prop:ig-frame-dependent-spectra-determinants}
\emph{(i) For $T=cI_{n}$ with $c>0$, which is a legitimate reframing since $\det(cI_K)=c^{K}>0$, the ordinary spectrum obeys $\operatorname{spec}\left(T^{-\top}\Lambda T^{-1}\right)=c^{-2}\operatorname{spec}(\Lambda)$. (ii) Let $U\subset\R^{n}$ have orthonormal basis $B$, let $U'=TU$ have orthonormal basis $B'$, and let $\Lambda'=T^{-\top}\Lambda T^{-1}$. Then}
\begin{equation}
\log\det\left(B'^\top\Lambda'B'\right)
=\log\det\left(B^\top\Lambda B\right)-\log\det\left(B^\top T^\top TB\right),
\label{eq:ig-logdet-law}
\end{equation}
\emph{so a subspace log-determinant is frame independent if and only if $T$ restricted to $U$ has unit determinant in absolute value; every orthogonal $T$ satisfies this condition.}
\status{ESTABLISHED}

\noindent\emph{Proof.} Part (i) is immediate. For (ii), write $TB=B'S$ with $S$ invertible, so $B'=TBS^{-1}$; orthonormality of $B'$ gives $S^{-\top}\left(B^\top T^\top TB\right)S^{-1}=I$ and hence $\det(S)^2=\det\left(B^\top T^\top TB\right)$. Then $B'^\top\Lambda'B'=S^{-\top}B^\top T^\top T^{-\top}\Lambda T^{-1}TBS^{-1}=S^{-\top}\left(B^\top\Lambda B\right)S^{-1}$, and taking log determinants gives \eqref{eq:ig-logdet-law}. \hfill$\square$

\medskip
The consequence is a prohibition, and it is the reason this section exists. A reframing changes no law; it changes only the components in which the law is written. A quantity computed from those components that is not invariant under \eqref{eq:ig-congruence} is therefore a property of the chart. By \Cref{prop:ig-frame-dependent-spectra-determinants}(i), any threshold criterion applied to the absolute eigenvalues of $\Lambda$ can be made to accept or to reject any given direction by choosing $c$, so no such criterion states anything about the model. Any admissible criterion must be phrased in generalized-eigenvalue form, that is as a statement about a pencil in which both members transform by the same congruence. By \Cref{prop:ig-frame-dependent-spectra-determinants}(ii) the same warning applies to determinant criteria evaluated on a subspace: they are invariant under orthogonal reframings and not otherwise, so a construction that uses one must either declare its frame or carry the Jacobian term in \eqref{eq:ig-logdet-law} explicitly. \status{ESTABLISHED}

\subsection*{Verification}

\Cref{prop:ig-generalized-spectrum-invariance} is proved above and needs no support. The deterministic replacement check \texttt{CHK-GENERALIZED-SPECTRUM} measures the generalized spectrum over twelve fully specified block congruences and uses the moving ordinary spectrum as its negative control. It establishes numerical usability only for that protocol; congruence invariance is proved above. \status{NUMERICAL}

\section{What is not claimed}
\label{sec:ig-notclaimed}

The Fisher metric is used here as a declared choice of geometry on the recognition manifold, and the classical uniqueness results are not strong enough to make that choice forced in this setting. On a finite sample space the Fisher metric is, up to a positive scale, the unique Riemannian metric invariant under Markov morphisms; this is Chentsov's theorem \citep{Cencov1982}, with the extension to non-normalized measures due to \citet{Campbell1986}. \status{ESTABLISHED} The latent space here is $\R^{n}$, so that theorem does not apply directly. Extensions to general sample spaces exist and are established, but they carry their own hypotheses: \citet{Ay2015} and \citet{Le2017} obtain uniqueness under a strong-continuity assumption on the information metric, and the monograph treatment collects the required regularity \citep{AyJostLeSchwachhofer2017}. \status{ESTABLISHED} The development therefore adopts the Fisher metric because it is the Hessian of the exponential-family potential \eqref{eq:ig-fisher-hessian} and because that is the structure the calculations use, and it does not claim that no other invariant metric is available in the continuous setting. \status{HYPOTHESIS}

No Markov/data-processing monotonicity theorem is asserted for the deterministic
Galerkin aggregation map. \status{NOT-CLAIMED} The established general statement is that the Fisher
metric contracts under a Markov kernel: pushing a family forward through a
stochastic map cannot increase Fisher information
\citep{AmariNagaoka2000,AyJostLeSchwachhofer2017}. \status{ESTABLISHED} That
theorem does not cover the aggregation operation of the following chapters,
because by \Cref{prop:ig-pullback-vs-pushforward} the coarse operator is a
restriction, a pullback along an inclusion of mean submodels, and not a
pushforward along a Markov kernel; the two differ by the positive semidefinite
Schur term in \eqref{eq:ig-restriction-marginal}, and the restriction is the
larger of the two in the Loewner order. \status{ESTABLISHED} Whether a monotone functional decreases
under the full declared aggregation/RG map is open. Closing it requires either
exhibiting a Markov kernel whose pushforward realizes the map, after which the
classical theorem applies verbatim, or exhibiting a Lyapunov functional and
proving its decrease directly. \status{OPEN} The later fixed-covariance Gaussian relaxation
in \Cref{sec:cg-gaussian-fixed-relaxation} is a different continuous-time
dynamics: its quadratic energy has an exact dissipation identity, and that
scoped flow is a Fisher natural-gradient flow. \status{ESTABLISHED} It does not turn the coarse map
itself into a Markov contraction. \status{NOT-CLAIMED}

Nothing here should be read as importing a uniqueness statement from the quantum setting. There, the metrics monotone under completely positive trace-preserving maps form a family parameterized by operator monotone functions rather than a single metric \citep{Petz1996}, so the phrase ``the monotone metric'' does not denote in that context, and an argument that relies on uniqueness cannot be transported from it. \status{ESTABLISHED} No Markov-morphism contraction theorem for the deterministic aggregation maps is asserted. \status{NOT-CLAIMED}

Finally, this section supplies a metric and not a dynamics. No natural-gradient
identification should be inferred merely from the appearance of \(\Lambda\) as
a preconditioner: an update is a natural gradient only when its metric is the
Fisher metric of the family being updated, in a named chart, and with the
objective differential typed in the matching chart
\citep{Amari1998,martens2020new}. The fixed-covariance Gaussian mean flow in
\Cref{sec:cg-gaussian-fixed-relaxation} satisfies those requirements and is
identified there. \status{ESTABLISHED} No broader identification for the remaining updates is
claimed. \status{NOT-CLAIMED}
